\documentclass[a4paper,11pt]{article}
\usepackage[utf8]{inputenc}
\usepackage[T1]{fontenc}
\usepackage{lmodern,amsmath,amssymb,amsthm,geometry,hyperref}
\geometry{margin=24mm}
\hypersetup{hypertexnames=false,colorlinks=true,urlcolor=blue,
pdftitle={Degeneracy, invariance, and exchange: proposed additions to v0.5}}
\setlength{\emergencystretch}{2em}
\newtheorem{proposition}{Proposition}
\newcommand{\norm}[1]{\left\lVert #1\right\rVert}
\newcommand{\ran}{\operatorname{ran}}
\newcommand{\spec}{\operatorname{spec}}
\begin{document}
\begin{flushleft}
Not incorporated into the delivered manuscript v0.5; not externally peer reviewed}
\vspace{0.7em}

{\Large\bfseries Degeneracy, invariance, and exchange: proposed additions to v0.5}
\end{flushleft}

\paragraph{Editorial purpose and status.}
When does a geometrically degenerate triplet sector support closed,
symmetry-compatible interaction dynamics?
A threefold linear eigenvalue, an isotropic projected coupling, and a mode
with low radiation are distinct properties.
The existing counterexamples in v0.5 can make this distinction the central
contribution. The two propositions below are finite-dimensional linear
algebra, with no novelty claim. They introduce no new numerical results.
The effective two-mode example in the final section motivates a future
question; it is not a derived model of the project's Q-balls.

The source proposal is \texttt{BEITRAG.txt}, sections A--C, supplied by
ag-phy-coordination on 30 September 2026. An independent algebraic reading
within OpenAI is documented in \texttt{ALGEBRA-REVIEW.txt}; it is not an
external scientific review. The original v0.5 manuscript and its artifacts
remain unchanged.

\section*{1. Eliminated modes and the symmetry of a projected sector}
Let \(H=H^\dagger\) be a time-independent operator on a finite-dimensional
complex Hilbert space. Let \(P=P^\dagger=P^2\) be an orthogonal rank-three
projector, and \(Q=I-P\). The projectors select subspaces; \(P\) is not
assumed to be a spectral projector of the full coupled \(H\).
In orthonormal bases of \(\ran P\) and \(\ran Q\), write
\begin{equation}
H=\begin{pmatrix}\lambda I_3&V\\ V^\dagger&D\end{pmatrix},
\qquad
\lambda\in\mathbb R,\quad D=D^\dagger,\quad V=PHQ.
\tag{S1}
\end{equation}
Here and below, block operators are restricted to their respective
subspaces. Norms are Euclidean vector norms and their induced operator norms.

\begin{proposition}[Schur correction and sufficient symmetry protection]
For real \(z\notin\spec(D)\), define
\begin{equation}
H_{\mathrm{eff}}(z)=\lambda I_3-V(D-zI_Q)^{-1}V^\dagger .
\tag{S2}
\end{equation}
Then
\begin{equation}
\det(H-zI)=\det(D-zI_Q)
           \det\!\left(H_{\mathrm{eff}}(z)-zI_3\right).
\tag{S3}
\end{equation}
Thus, away from the excluded poles, the full eigenvalue problem is
equivalent to the reduced eigenvalue condition.
Writing \(\delta=\operatorname{dist}(z,\spec(D))>0\), one has
\begin{equation}
\norm{H_{\mathrm{eff}}(z)-\lambda I_3}
\leq \frac{\norm{V}^{\,2}}{\delta}.
\tag{S4}
\end{equation}
If a unitary group representation \(R(g)\) satisfies
\([R(g),H]=[R(g),P]=0\) for every \(g\), and its restriction to
\(\ran P\) is irreducible over \(\mathbb C\), then
\(H_{\mathrm{eff}}(z)\) is scalar on \(\ran P\).
\end{proposition}

\begin{proof}
For a full eigenvector \((u,v)\), the lower block equation gives
\(v=-(D-zI_Q)^{-1}V^\dagger u\). Substitution into the upper block yields
\((H_{\mathrm{eff}}(z)-zI_3)u=0\). At the allowed values of \(z\), a
nonzero full eigenvector cannot have \(u=0\), and each reduced eigenvector
reconstructs a full one by this formula. Block elimination gives (S3).
Since \(D\) is Hermitian, the spectral theorem gives
\(\norm{(D-zI_Q)^{-1}}=1/\delta\), proving (S4).

Commutation with \(P\) makes \(R(g)\) block diagonal,
with restrictions \(R_P(g)\) and \(R_Q(g)\). Commutation with \(H\)
implies \(R_PV=VR_Q\) and \([R_Q,D]=0\).
Consequently \(V(D-zI_Q)^{-1}V^\dagger\) commutes with every \(R_P(g)\).
Schur's lemma on the irreducible complex triplet then makes this
operator, and hence \(H_{\mathrm{eff}}(z)\), scalar.
\end{proof}

\paragraph{Explicit counterexample to protection by projection alone.}
Take a single complementary mode \(D=\Delta\), with
\(\Delta\in\mathbb R\), and \(V=g(1,0,0)^T\), \(g\in\mathbb C\setminus\{0\}\).
For \(z\ne\Delta\),
\begin{equation}
H_{\mathrm{eff}}(z)-\lambda I_3
=-\frac{|g|^2}{\Delta-z}
  \begin{pmatrix}1&0&0\\0&0&0\\0&0&0\end{pmatrix}.
\tag{S5}
\end{equation}
One triplet direction is distinguished even though the direct compression
\(PHP\) is exactly isotropic. For real \(g\), the numerator is \(g^2\).
Equal eigenvalues of \(PHP\) therefore do not by themselves protect the
symmetry of the reduced coupled problem.

\paragraph{Limits of the statement.}
The determinant reduction is not to be evaluated at a pole \(z\in\spec(D)\).
The norm estimate controls the entire correction, not a compulsory amount
of anisotropy; a nonzero correction can be isotropic. A full symmetry
representation can protect isotropy while still allowing transfer to
equivalent representations inside \(Q\). It does not in itself imply
\(QHP=0\) or absence of leakage.
Moreover, \(H_{\mathrm{eff}}(z)\) is an energy-dependent Schur operator,
not automatically a time-local generator for the projected evolution.
The actual Q-ball frequency pencil and its normalization must be derived
and treated on their own terms. This finite Schr\"odinger-type formula
cannot be imported into that problem without justification.

\section*{2. A conservative finite-time leakage bound}
\begin{proposition}[Leakage from an initially occupied subspace]
For the same time-independent Hermitian \(H\), let
\(i\dot\psi=H\psi\), \(\norm{\psi(0)}=1\), and \(Q\psi(0)=0\).
Then, for every real \(t\),
\begin{equation}
Q\psi(t)=-i\int_0^t
 e^{-iD(t-s)}QHP\,P\psi(s)\,ds,
\tag{S6}
\end{equation}
and
\begin{equation}
\norm{Q\psi(t)}\leq
\min\{1,\,|t|\norm{QHP}\},\qquad
\norm{Q\psi(t)}^2\leq
\min\{1,\,t^2\norm{QHP}^{\,2}\}.
\tag{S7}
\end{equation}
These statements hold for any orthogonal projector \(P\), not only rank three.
\end{proposition}

\begin{proof}
Projecting the evolution gives
\(i\,d(Q\psi)/dt=DQ\psi+QHP\,P\psi\).
Variation of constants, with zero initial \(Q\)-component, yields (S6).
Both the full evolution and \(e^{-iD(t-s)}\) are unitary.
Thus \(\norm{P\psi(s)}\leq1\), and the integral has norm at most
\(|t|\norm{QHP}\). For negative \(t\), reverse the oriented integration
interval; its length is still \(|t|\).
The full-state normalization independently gives
\(\norm{Q\psi(t)}\leq1\). Combining the two bounds and squaring proves (S7).
\end{proof}

\paragraph{Connection to the existing invariance checks.}
For classical complex amplitudes, \(\norm{Q\psi(t)}^2\) is the fraction of
the normalized squared norm outside the chosen subspace. Calling it a
probability additionally presumes a quantum-state interpretation.
The exact condition \(QHP=0\) implies zero leakage for every initial state
in \(P\). For an observation window \(|t|\leq T\), \(T>0\), a sufficient
condition for an outside norm fraction at most \(\varepsilon\in[0,1]\) is
\(\norm{QHP}\leq\sqrt{\varepsilon}/T\).
Reported Frobenius norms of outside blocks may be used as conservative
upper bounds on the induced norm, provided the projector, normalization,
and block definition actually match.
This connects the algebraic invariance criterion to a finite tolerated
error. It proves neither sharp long-time stability nor an optimal
lifetime. For large times the bound may reduce to the trivial value one.
Nonlinear fields require other error estimates; no nonlinear or
continuum certificate follows from (S7).

\section*{3. A future question: dark and still interacting?}
As a purely effective passive two-mode illustration, consider
\begin{equation}
\mathcal H=
\begin{pmatrix}
\omega-i\gamma/2 & J-i\gamma_{12}/2\\
J-i\gamma_{12}/2 & \omega-i\gamma/2
\end{pmatrix},
\quad
\omega,J,\gamma,\gamma_{12}\in\mathbb R,\quad
\gamma\geq|\gamma_{12}|.
\tag{S8}
\end{equation}
The last condition makes the loss matrix positive semidefinite.
The normalized symmetric and antisymmetric vectors have eigenvalues
\begin{equation}
z_\pm=\omega\pm J-\frac{i}{2}(\gamma\pm\gamma_{12}).
\tag{S9}
\end{equation}
For \(\gamma_{12}=\gamma\), the antisymmetric channel has zero decay
in this model even when \(J\ne0\). Thus zero decay in a channel is not
logically equivalent to absent coherent coupling. This is a familiar
algebraic possibility, not a parameter choice derived from KRAFT or
from the project's microscopic field model.

In a microscopic mediation mechanism, the dispersive and dissipative
parts arise from the same causal response and cannot simply be chosen
independently. Nor does a fixed \(J\) establish a spatial force: that
requires positional dependence, dynamical positions, and consistent
energy and momentum accounting. The coordinator reports that
KRAFT-still failed its strict force-switch criterion; this example
does not reinterpret that result.

The smallest proposed next question is whether the same field or
scattering operator yields both finite coherent frequency splitting and
small radiation, with a consistently defined distance force where
applicable. Frequency splitting, radiation channels, and energy flux
must refer to the same configuration. Different phase windows or fits
must not be presented as one joint finding. This is a proposed study,
not a calculation performed for this supplement.

\section*{4. Placement and boundaries}
Propositions 1 and 2 are proposed for insertion after the exchange and
subspace-invariance criteria of v0.5, section 4. They formalize why a
degenerate compression or an isotropic projected matrix is insufficient
when complementary modes remain coupled.
Section 3 belongs among the concrete future field-model questions,
not among the already completed triplet experiments.

Frequency-dependent mediator range, Hamiltonian phase memory, and
Hawking-related questions are retained in the separate, attributed
\texttt{RESEARCH-NOTE-EN.txt}. They do not extend the claims of the main
paper. No KRAFT numbers are mixed into the triplet results, and the
existing BIC/exterior proof remains a separate line of work.
\end{document}
